% #############################################################################
% This is Chapter 9 - Evaluation
% !TEX root = ../main.tex
% #############################################################################
\fancychapter{Evaluation}
\label{chap:evaluation}

This chapter reports the Evaluation activity of the \ac{DSRM}. Following Pries-Heje et al.'s distinction between \emph{when} and \emph{how} an artefact is evaluated~\cite{J.Pries-Heje:2008SD}, the evaluation carried out here is \emph{ex post}, since it assesses the artefact after construction and on the basis of its actual use, and \emph{naturalistic}, since it takes place in a real organisational setting with real users rather than in a laboratory.

% #############################################################################
\section{Evaluation Strategy}
\label{sec:eval_strategy}

Two artefacts were produced in \Cref{chap:proposal,chap:apex}: a framework and its instantiation. They are not evaluated by the same means, and conflating them would weaken both claims. \Cref{tab:eval_instruments} therefore states which instrument evaluates which artefact, and which research question it addresses.

{\small
\setlength{\tabcolsep}{5pt}
\begin{longtable}{>{\raggedright\arraybackslash}p{4.4cm}
                  >{\raggedright\arraybackslash}p{3.1cm}
                  >{\raggedright\arraybackslash}p{1.6cm}
                  >{\raggedright\arraybackslash}p{5.4cm}}
\caption{Evaluation instruments, target artefact, and research question addressed}
\label{tab:eval_instruments} \\
\toprule
\textbf{Instrument} & \textbf{Artefact evaluated} & \textbf{DRQ} & \textbf{What it establishes} \\
\midrule
\endfirsthead
\toprule
\textbf{Instrument} & \textbf{Artefact evaluated} & \textbf{DRQ} & \textbf{What it establishes} \\
\midrule
\endhead
\midrule
\multicolumn{4}{r}{\textit{Continued on next page}}\\
\bottomrule
\endfoot
\bottomrule
\endlastfoot

Demonstration in a real project & Framework and instantiation & DRQ3, DRQ4 & That the artefact is applicable and produces the intended artefacts and decisions \\
\addlinespace
\ac{SUS} & Instantiation (Apex) & DRQ4 & Perceived usability of the supporting tool \\
\addlinespace
\ac{NASA-TLX} & Instantiation (Apex) & DRQ4 & Perceived cognitive workload imposed on the practitioner \\
\addlinespace
Apex \ac{UX} questionnaire (author-designed) & Instantiation (Apex), with two items scoped to the framework & DRQ4 & Trust in AI-generated content, legibility of phase gates, and error surfacing, which \ac{SUS} and \ac{NASA-TLX} are not designed to reach \\
\addlinespace
Practitioner interviews & Framework & DRQ1, DRQ2 & Practicality, feasibility, and perceived value in real workflows \\
\addlinespace
Agile and Scrum expert interviews & Framework & DRQ1, DRQ2 & Alignment with, and departure from, established delivery methodology \\
\addlinespace
Analytical assessment against criteria & Framework & DRQ1, DRQ2 & Degree to which stated design criteria are satisfied \\
\end{longtable}
}

% This distinction is worth defending explicitly in the text, because it is the
% obvious line of attack: a good SUS score for the tool is NOT evidence that
% the framework is sound. Say so before an examiner does.

\ac{SUS} and \ac{NASA-TLX} measure Apex, not the framework. Both are administered to someone who has spent a session operating a user interface, and neither contains an item about phase decomposition, gate placement, or accountability for generated artefacts. A favourable \ac{SUS} score therefore establishes that the tool is usable and nothing at all about whether the method it implements is sound; an unfavourable one, correspondingly, is not evidence against the method. That is a design decision rather than a defect: the framework's own validity is carried by the practitioner and expert interviews (\Cref{sec:interviews}), the analytical assessment against criteria fixed in advance (\Cref{sec:analytical}), and the demonstration (\Cref{chap:demonstration}), each of which reaches the method directly.

One instrument spans both artefacts, and the split is handled rather than ignored: two items of the author-designed \ac{UX} questionnaire ask about the process rather than the product, and are analysed apart from the fifteen instantiation-scoped ones, as \Cref{sec:apex_ux} sets out.

% #############################################################################
\section{Participants}
\label{sec:eval_participants}

Participants were recruited through the supervisor's own professional connections. There was no open call, no advertisement and no mailing list, and the number of people approached was not recorded, so no response rate can be reported and none is claimed. Thirteen sessions were completed, coded P0 to P12; all thirteen returned all four instruments and no participant withdrew. P0 was planned as an uncounted pilot and is counted as a full participant, because the session produced complete, correctly labelled data under the same consent form, task script and instruments every later participant received, and discarding usable data on the strength of its position in the sequence alone had no principled basis. The pre-registered target was twelve; the achieved sample is thirteen, and the sample is not drawn from a single organisation but from one person's professional network, which is a different limit rather than a weaker one.

The sample is junior and mixed. On years of professional experience, three participants were students, three had under one year, and seven had one to three years; nobody had more than three. By current role, six were students, four developers and one an \ac{AI} engineer, and two came from outside software development entirely, one a communications network technician and one a partner-manager in a hairdressing business. Prior exposure to AI coding assistants, pre-registered as the strongest expected confound, spread across the whole scale: four used them daily, four weekly, two monthly, two had tried one once or twice, and one had never used one. Familiarity with the notations and tools the session assumes was low by design of the sample rather than by selection: nine of the thirteen had no familiarity with Gherkin or \ac{BDD} and four had seen it, with nobody reporting regular use; ten reported no familiarity with Taiga, the project-management tool the session runs against, two some, and one left the item blank. Nine participants chose English as their session language and four Portuguese, and ten completed the instruments in English against three in Portuguese; the two participants for whom those two facts disagree are reported in \Cref{sec:threats} rather than reconciled here.

What that sample supports is bounded. Thirteen respondents meets the twelve at which a \ac{SUS} mean is generally treated as stable, so the mean is reported rather than suppressed, with the distribution and every individual score alongside it. It remains a convenience sample, skewed junior and assembled from a single network, and what that does and does not support is set out in full in \Cref{sec:threats}; no inferential statistic is computed anywhere in this chapter and no subgroup is compared against another.

Every participant gave informed consent before any data was collected, as three
separate tick boxes on the first page of an online form completed before the task
script was opened, so that no response was collected ahead of consent and
agreeing to take part did not imply agreeing to be quoted or to be contacted for
a follow-up interview. The form is reproduced in \Cref{chapter:appendixB}. All
thirteen consented to participate; eleven consented to anonymised quotation and
eleven to being contacted, the two sets overlapping at nine, and quotations below
are drawn only from the eleven. No submission to an institutional ethics
committee was required: the supervisor confirmed that for a study of this kind
informed consent, together with anonymisation and a stated retention period, is
sufficient at \ac{IST}, and that no standard institutional consent form exists
that would supersede the one used here.

No session was observed and nothing was recorded: no screen capture, no audio
and no video, of the task sessions or of the interviews, and the consent form
contains no recording item because there was nothing to seek permission for.
What was collected is the questionnaire responses, the timestamps
at which each form was submitted, and, for the interviews, contemporaneous
written notes taken by the researcher and written up in full immediately after
each interview. The methodological cost of this choice, principally that a set of
written notes is selective in a way a full transcript is not, is set out in
\Cref{sec:threats}.

Participants are identified only by a code of the form P1 to Pn. Names,
employers, and identifying project details were not recorded with the responses,
and quotations reproduced in this chapter are anonymised, with any detail that
could identify a participant or their organisation removed or generalised. Only
wording written down verbatim at the time is presented as a quotation; the
remainder of the interview material is presented as paraphrase and is marked as
such. The anonymised questionnaire responses and interview notes are retained,
accessible only to the researcher, and are available as stated in
\Cref{sec:appB_raw_scores}, so that the aggregate figures reported in
\Cref{sec:eval_results} can be independently recomputed.

% #############################################################################
\section{Instruments and Procedure}
\label{sec:eval_instruments_procedure}

The usability study was administered entirely unmoderated. Each participant
worked alone, at their own pace, from a fixed written task script, and submitted
the questionnaires as online forms without the researcher present. A second,
moderated arm of a small number of observed think aloud sessions was designed and
deliberately dropped, for two reasons. First, no session was to be
recorded, and an observed session without a recording yields notes written by a
facilitator who is simultaneously facilitating, which is the weakest form of the
only evidence such an arm exists to produce. Second, and more seriously, the
author designed both artefacts under evaluation; an unmoderated design removes
that conflict structurally rather than mitigating it, since no facilitator is
present to steer, hint, or react. What the decision costs, namely
measured task completion, time on task, assist counts, and any observational
account of where participants struggled, is stated in \Cref{sec:threats}.

\subsection{System Usability Scale}
\label{sec:sus}

The \ac{SUS}~\cite{Brooke:1996SUS} is a ten-item instrument scored on a five-point scale of strength of agreement, administered after the participant has completed a defined set of tasks. Participants responding in Portuguese are given the validated European Portuguese version~\cite{Martins:2015SUS} and those responding in English the original wording, both as reproduced in \Cref{sec:appB_sus}; the two language versions are reported separately rather than pooled.

Each item contributes between zero and four points: for the odd-numbered, positively worded items the contribution is the response minus one, and for the even-numbered, negatively worded items it is five minus the response. The ten contributions are summed and multiplied by 2.5, giving a value between 0 and 100. That value is not a percentage: 100 is not a perfect result and 50 is not half of anything, which is why it is not left to speak for itself.

Two published interpretation scales are applied to it, because they say different things about the same figure and reporting only one invites the more flattering reading. Bangor et al. attach adjectives to score ranges from a corpus of over two thousand administrations, placing \emph{Poor} at 35.7, \emph{OK} at 50.9, \emph{Good} at 71.4 and \emph{Excellent} at 85.5~\cite{Bangor:2009SUS}. Sauro and Lewis convert a score into a percentile against their own corpus and assign a curved grade from A to F, in which the corpus mean of approximately 68 sits at the fiftieth percentile rather than at a pass mark~\cite{Sauro:2016QUANT}. Both are reported in \Cref{sec:eval_results}.

The two-factor decomposition of Lewis and Sauro is reported alongside the overall score: Usability from the eight items other than 4 and 10, rescaled by 100/32, and Learnability from items 4 and 10, rescaled by 100/8~\cite{Lewis:2009SUS}. It earns the space because the two sub-scales can move in opposite directions, and a tool that is hard to learn but workable once learnt is a different finding, with a different remedy, from one that is neither.

The Portuguese results carry a caveat the validation study itself states. Construct validity was established against the Post-Study System Usability Questionnaire at $r = 0.70$, but inter-rater reliability was weak, with an intraclass correlation coefficient of 0.36 and a 95 per cent confidence interval of 0.01 to 0.63, and percentage agreement reached 76.67 per cent against the authors' own 80 per cent threshold, which they attribute to the alternating polarity of the items; the validation sample was drawn from the general community rather than from software practitioners~\cite{Martins:2015SUS}. With only three Portuguese respondents in this study, the property that instrument is weakest on is exactly the one a three-person mean depends on most, and the Portuguese figures are accordingly reported separately and read as indicative only.

\subsection{NASA Task Load Index}
\label{sec:tlx}

The \ac{NASA-TLX}~\cite{Hart:1988TLX} measures perceived workload across six subscales: mental demand, physical demand, temporal demand, performance, effort and frustration. This evaluation uses the unweighted Raw \ac{NASA-TLX} variant~\cite{Hart:2006TLX}, omitting the fifteen pairwise weighting comparisons. The instrument is administered once after each of the marked tasks rather than once at the end of the session, which is the instrument's intended per-task use and yields a workload profile per lifecycle phase rather than a single aggregate.

The Raw variant follows from that choice of administration: those fifteen comparisons would become ninety per participant under per-task administration, well beyond what the session can absorb, and the weights earn their cost primarily when workload is compared across experimental conditions, which this evaluation has none of. It is not, however, adopted only on grounds of feasibility. Reviewing five hundred and fifty studies twenty years after the instrument's publication, Hart records that removing the weighting is the most common modification made to the \ac{NASA-TLX}, and that across the twenty-nine studies in which the two versions were directly compared the Raw variant was found to be more sensitive, less sensitive, or equally sensitive depending on the study, with no consistent advantage to weighting~\cite{Hart:2006TLX}. The subscale scores, which are this evaluation's primary result, are in any case identical under both versions, since the weighting affects only how they are combined into a single aggregate.

Two departures from the published instrument are recorded here. Workload was collected on a linear scale from 0 to 10 rather than on the original twenty-interval line, and the responses were rescaled to the instrument's 0 to 100 range by multiplying by ten; no form tool suitable for an unmoderated study reproduces a twenty-interval line. The subscales were also reordered to place the reversed performance scale last rather than fourth, because a reversed scale encountered in the middle of five forward ones, with no facilitator present to explain it, is the most plausible way a participant inverts a sixth of the data set. Neither change affects the unweighted Raw \ac{NASA-TLX} computation. The first, however, reduces the resolution of the instrument: every response is a multiple of ten, so the resulting values are not directly comparable to studies using the original response scale, and small differences between cells of the workload matrix are not interpreted as meaningful.

Repeated administration carries its own cost. Rating the same six dimensions six times invites anchoring, where later ratings drift towards earlier ones, and fatigue on the final tasks of the session; neither is avoidable under per-task administration, and both are carried into \Cref{sec:threats} and weighed there against the workload profile that administration is what makes obtainable.

Both instruments are administered against one fixed task set, identical in wording and order for every participant. Nine tasks carry a single Unit of Work, a password-reset feature, from a written requirement through to a deployed increment: sign in and open the seeded project (1), connect it to a GitHub repository (2), take the requirement to locked Gherkin (3), produce and lock the design (4), sign out and resume the session (5), break the story into implementation tasks (6), generate a test plan and record the \ac{QA} sign-off (7), deploy the feature (8), and export the resulting artefacts (9). Six of the nine, tasks 3, 4, 6, 7, 8 and 9, each carry a Raw \ac{NASA-TLX} form; tasks 1, 2 and 5 are setup and session-resumption steps and carry none. All seventy-eight administrations the design calls for, thirteen participants by six marked tasks, were received, so no workload cell below rests on a partial set of respondents. The project environment was reset to an identical seeded state before every session, so that no participant met another's work.

One change to the task set is recorded here because it removed a measure instead of relocating one. Task 8 was originally a refusal test: the participant was to attempt to deploy a story that had deliberately never passed \ac{QA}, and be blocked at the gate. It was redesigned on 30 August 2026 so that tasks 3 to 9 form one continuous flow through the same feature, which means the deployment in task 8 succeeds, and succeeds \emph{because} of the \ac{QA} sign-off the participant gave in task 7; the link between those two gates is what the redesigned task tests. What the redesign cost the \ac{UX} questionnaire is stated in \Cref{sec:apex_ux} and its consequence for the study is carried into \Cref{sec:threats}.

\subsection{Apex UX Questionnaire}
\label{sec:apex_ux}

An eighteen-item questionnaire was written for this evaluation to reach what
\ac{SUS} and Raw \ac{NASA-TLX} are not designed to say: whether a participant
trusted AI-generated content, whether a phase gate read as a guardrail or as
obstruction, and whether reviewing generated work felt cheaper than writing
it. Its items are grouped into five constructs: understanding where you are
in the process, working with AI-generated content, control and
reversibility, feedback and errors, and fit to real work. It is not a
validated instrument and is not treated as one: it is reported item by item,
with no composite score or subscale means, since collapsing it into a
single figure would claim more precision than it supports. Administered last, after \ac{SUS}, and
reproduced in full in \Cref{chapter:appendixB}.

Two items in the fifth construct, on whether the artefacts and the process
fit real work, ask about the framework rather than the instantiation and
are analysed under the framework's evidence rather than Apex's
(\Cref{tab:eval_instruments}): keeping them on the same sheet is
convenient for the participant, splitting them out in analysis is what
keeps the instrument from quietly undermining that separation.

One item pair is expected in advance to be the most diagnostic in the instrument. Item B4 asks whether the participant trusted the generated content enough to use it in a real project, and item B2 whether reviewing that output took less effort than writing it. Neither is interesting alone; read together they are a trust-calibration pair. High trust reported alongside low review effort is the signature of acceptance without scrutiny, which is the failure mode the framework's gates exist to prevent and which no participant will report directly when asked. High trust alongside high review effort is the opposite and far less worrying reading: content believed because it was checked. The pair is the closest an unmoderated, self-report design can get to the observation that task 4 was designed to elicit, namely whether generated output is edited or accepted wholesale.

A second cross-read was designed into the instrument and is no longer available. Item A4, \emph{when the tool refused to let me continue, I understood why}, was to be read against the \ac{NASA-TLX} Frustration score and free-text response for task 8, which was at that point a deliberate refusal test; with no observer present and no scored comprehension check, that triangulation was the study's only evidence on whether a phase gate reads as a guardrail or as obstruction. When task 8 was redesigned (\Cref{sec:tlx}) the guaranteed refusal went with it, and A4 was removed from both language versions on 30 August 2026 rather than left to be answered by participants who had not been refused anything. Seventeen of the eighteen items written were therefore administered, and A4 is not missing data, since the item does not exist for any of the thirteen participants. What was lost is not one item but the study's only measure of gate comprehension, and that loss is carried into \Cref{sec:threats}.

\subsection{Semi-Structured Interviews}
\label{sec:interviews}

Interviews were conducted with a population separate from the unmoderated task-based study reported in \Cref{sec:eval_participants}: practitioners and an Agile/Scrum expert recruited individually by the researcher rather than as a follow-up to that arm. This is a deliberate design choice rather than a gap. The interview guides in \Cref{chapter:appendixB} ask about the framework's practicality, its adoption cost, and its alignment with established delivery methodology; these are questions a participant who has just spent an hour on nine fixed tasks is poorly placed to answer for their own team and organisation, and the Agile-expert guide explicitly does not require the tool to have been used at all. The themes below were derived from the written notes described in \Cref{sec:eval_participants}: the codebook was developed from the first write-ups in each group and applied across the remainder, with codes added as they emerged. Session duration was not logged systematically and is not reported.

For the same anonymity reason stated in \Cref{sec:eval_participants}, interviewees are identified here only by a code: T1 to T5 for practitioners and E1 for the Agile/Scrum expert. These are kept deliberately separate from the P-codes used for the unmoderated study's participants, since the two populations are disjoint and neither did the other's protocol.

Five practitioners and one Agile expert had been interviewed at the time of writing, against Guide 1's target of five to eight practitioners and Guide 2's target of two to three experts; one practitioner, T3, was interviewed a second time after the framework had been substantially revised in response to the first round. Both groups therefore sit at or below the low end of their target, and recruitment continues past submission of this draft; the sample is reported at its current size rather than held back for a target that would postdate the deadline. What follows should be read at that scale, as a small number of attributable, individually rich sessions rather than a saturated qualitative sample, and the themes below name which interviewee raised which point precisely because the sample is too small for that attribution to be smoothed away by aggregation. Two of the five practitioners, T1 and T5, are not quoted below: their sessions took place earlier, during the tool's construction, and their material is reported instead in \Cref{sec:apex_history}, where it did concrete work in changing the design. T2, T3 and T4 are the three whose interviews are this section's evidence.

\paragraph{Practitioners.} T2 offered the framework's central bet, that a locked specification should be the durable artefact and code should be regenerable from it, its clearest unprompted endorsement in this group: specifications, in his description, function like a library for project information and as the vehicle for guidelines, requirements, plans and tests kept together. The same practitioner separately named a real limit on that same idea, that it is difficult to account for every assumption a specification leaves implicit, and suggested grounding requirements by asking the user a structured series of questions to establish that foundation before generation begins.

T4 raised semantic versioning of the specification as a requirement rather than a suggestion, arguing that a human must always be able to consult and edit an artefact's history rather than trust it only as it currently stands. The same practitioner named, unprompted, the risk that a strictly linear phase sequence reintroduces waterfall-style rigidity, arguing for feedback loops rather than a one-directional pipeline. This second point is worth isolating because the Agile-expert guide in \Cref{chapter:appendixB} was written anticipating exactly this criticism from methodology experts specifically; a practitioner raising it unprompted, without an expert being asked to probe for it, is stronger evidence that the concern is substantive than if it had only ever come from a session designed to elicit it.

T3's first interview raised two further points, feature-level rather than framework-level and each reported for completeness rather than as a developed theme: that AI-generated requirements need more detail specific to the project-management tool in use, Taiga in this case, and that the design should incorporate more UI factors to keep on-screen state clearly defined. T3's second interview produced the sharpest challenge the framework received in this study, and its content, together with the design changes it is directly traceable into, is reported in full in \Cref{sec:apex_history} rather than here.

\paragraph{Agile/Scrum expert.} The one expert interview obtained, following a live presentation of the framework and a demonstration of Apex, produced the clearest validation this study collected of the demonstration argument made independently in \Cref{chap:demonstration}: that a tool instantiating and applying a framework's roles and phases makes the framework legible in a way a description on paper does not. The specific comparison volunteered was against the rigid, closed-role structures typical of methodology adopted inside large organisations, against which the framework's hats and its refusal to lock roles to named individuals read as a genuine point of difference; the Mob Elaboration ritual and the framework's governance metrics drew the same reaction, read as departures from what the expert's own organisation does today.

Two further observations from the same session are reported as open limitations rather than folded into that validation, because neither produced a design change. The first concerns Apex directly: navigating the tool mid-demonstration produced visible uncertainty about where a given action would lead next, and the expert named tooltips or contextual framing, and a tutorial, as what was missing. This is a first-contact observation from a single session rather than a defect evidenced the way \Cref{sec:apex_limitations} otherwise requires, and it is recorded as a gap rather than a measured finding. The second concerns the expert's own organisation rather than Apex: AI tool use there is centrally governed, restricted to a single permitted tool, and metered by token consumption, yet no role or specialist owns that governance, a gap large enough that the organisation is actively hiring against it. The expert raised, unprompted, whether AI participation in a lifecycle nets out as a cost or a benefit; the framework's governance metrics (\Cref{sec:metrics}) are positioned to answer exactly that question but do not yet operationalise it for token cost specifically, and the point is carried forward as a direction for future work in \Cref{chap:conclusion} rather than pursued further here.

\paragraph{Read across both groups.} The strongest convergent signal is the linearity concern: raised unprompted by a practitioner using the tool, and anticipated independently in the design of the expert guide before any interview took place, it reaches this study from two directions that were never told to look for the same thing. The strongest divergence spans two sections rather than sitting inside this one: T2's unprompted endorsement of the specification-as-source-of-truth construct is reported above, while the sharpest doubt this study collected on the same construct, from the same population, is reported in \Cref{sec:apex_history}. Naming that split matters, because the favourable reading is not reported as though it were uncontested.

\subsection{Analytical Assessment}
\label{sec:analytical}
The framework is additionally assessed against five criteria fixed before the results were seen rather than derived from them. \emph{Clarity} asks whether each construct is defined precisely enough to be applied without its author present. \emph{Completeness} asks whether the framework covers the lifecycle it claims to cover, including its own defined measurements. \emph{Alignment with \ac{SDLC} phases} asks whether the phase decomposition maps onto recognised lifecycle stages without merging, splitting or omitting one. \emph{Support for human-in-the-loop practice} asks whether the framework makes human review of generated output structurally necessary rather than merely advisable. \emph{Governance of AI-enabled activity} asks whether AI contributions are made traceable, attributable and measurable by mechanism rather than by discipline.

Each is rated on a three-level ordinal scale. \emph{Met} means the criterion is satisfied and evidence exists that it is. \emph{Partially met} means it is satisfied in part and the shortfall is identifiable rather than merely suspected. \emph{Not met} means it is not satisfied. No finer scale is used, because the evidence available cannot distinguish more than three levels and a five-point scale would imply a precision this assessment does not have.

Each rating is justified against evidence that exists independently of the author's opinion of his own design: the construct-to-mechanism mapping in \Cref{tab:apex_mapping}, whose negative rows are as usable here as its positive ones; the interviews in \Cref{sec:interviews}; and incidents recorded during the study sessions themselves. Where the only available support for a rating would be the author's own assessment, the rating is lowered rather than argued, on the principle that a criterion which cannot be evidenced is not distinguishable from one that is not met. The ratings are reported in \Cref{tab:analytical}.

% #############################################################################
\section{Results}
\label{sec:eval_results}

\subsection{System Usability Scale}
\label{sec:results_sus}

Thirteen participants returned complete \ac{SUS} responses with no blank item. \Cref{tab:sus_scores} gives every score individually, since at this sample size the distribution carries more information than the mean, which alone would conceal how wide it is.

{\small
\renewcommand{\arraystretch}{1.2}
\setlength{\tabcolsep}{6pt}
\begin{longtable}{l c r r r}
\caption[\acs{SUS} score and sub-scales, per participant]{\acs{SUS} score and the Usability and Learnability sub-scales, per participant, by language of administration. All three quantities run from 0 to 100 and none is a percentage.}
\label{tab:sus_scores} \\
\toprule
\textbf{Participant} & \textbf{Language} & \textbf{\acs{SUS}} & \textbf{Usability} & \textbf{Learnability} \\
\midrule
\endfirsthead
\toprule
\textbf{Participant} & \textbf{Language} & \textbf{\acs{SUS}} & \textbf{Usability} & \textbf{Learnability} \\
\midrule
\endhead
\midrule
\multicolumn{5}{r}{\textit{Continued on next page}}\\
\bottomrule
\endfoot
\bottomrule
\endlastfoot

P0 & EN & 42.5 & 50.00 & 12.5 \\*
P1 & EN & 75.0 & 81.25 & 50.0 \\*
P5 & EN & 80.0 & 87.50 & 50.0 \\*
P6 & EN & 15.0 & 12.50 & 25.0 \\*
P7 & EN & 62.5 & 71.88 & 25.0 \\*
P8 & EN & 60.0 & 59.38 & 62.5 \\*
P9 & EN & 60.0 & 65.63 & 37.5 \\*
P10 & EN & 75.0 & 75.00 & 75.0 \\*
P11 & EN & 65.0 & 56.25 & 100.0 \\*
P12 & EN & 70.0 & 75.00 & 50.0 \\ \addlinespace
P2 & PT & 37.5 & 46.88 & 0.0 \\*
P3 & PT & 40.0 & 46.88 & 12.5 \\*
P4 & PT & 55.0 & 50.00 & 75.0 \\ \addlinespace
\midrule
\textbf{Mean, all 13} & & \textbf{56.73} & \textbf{59.86} & \textbf{44.23} \\*
Mean, English ($n = 10$) & & 60.50 & & \\*
Mean, Portuguese ($n = 3$) & & 44.17 & & \\
\end{longtable}
}

The mean across all thirteen is \textbf{56.73}, with a standard deviation of 18.50, a median of 60 and a range from 15 to 80. That is below the published corpus mean of approximately 68: on the curved grading scale it is a \textbf{grade D}, in the fifteenth to thirty-fourth percentile band, and on the adjective scale it sits between \emph{OK} at 50.9 and \emph{Good} at 71.4, closer to the former. The two language groups differ, and are not pooled for interpretation: the ten English responses average 60.50 (SD 19.18, median 63.75), also grade D, while the three Portuguese responses average 44.17 (SD 9.46, median 40), which is a \textbf{grade F}, in the zeroth to fourteenth percentile band. Three respondents cannot support a claim about Portuguese-speaking users, particularly given the agreement limits that instrument's own validation reports, but the figure is stated rather than withheld.

The spread is the more informative part: six of the thirteen scores are above 60, two are at 60, and five are below it, with four at or above 70, the region the adjective scale calls \emph{Good}. One participant, P6, scored \textbf{15}, below Bangor's \emph{Awful} anchor at 20.3 and near \emph{Worst imaginable} at 12.5. That score is reported as recorded and is not excluded, and it should be read alongside one observation about the same participant's other data: of the thirty-six Raw \ac{NASA-TLX} ratings P6 gave, thirty-five are the identical value, which is the signature of a constant response style rather than a differentiated judgement. The observation cuts both ways and neither is pursued, since thirteen participants give no capacity to absorb one atypical respondent and no grounds to remove one.

The sub-scales separate in the way the decomposition exists to expose. Usability averages 59.86 and Learnability 44.23, so the weaker of the two is the one concerning how much has to be learned before the tool can be used at all. Learnability is also far more dispersed, running the full 0 to 100: P2 scored 0, maximally agreeing both that they would need technical support to use the system and that they would need to learn a great deal before getting going, while P11 scored 100, disagreeing with both as strongly as the scale allows. Usability by comparison runs from 12.50 to 87.50. The reading this supports, and which \Cref{sec:eval_discussion} takes up against the other instruments, is that the difficulty reported is concentrated at first contact rather than distributed evenly through use.

\subsection{Raw NASA Task Load Index}
\label{sec:results_tlx}

\Cref{tab:tlx_matrix} is the primary workload result: the six subscales against the six marked tasks, each cell a mean over all thirteen participants. Reporting the aggregate alone would discard the part of the instrument that does the work here, since the six subscales move independently and their shape is the finding.

{\small
\renewcommand{\arraystretch}{1.2}
\setlength{\tabcolsep}{5pt}
\begin{longtable}{l r r r r r r r}
\caption[Raw \acs{NASA-TLX} subscale means by task]{Raw \acs{NASA-TLX} subscale means by task, over thirteen participants, on the instrument's 0 to 100 range. Performance carries reversed anchors and is reported as marked, so a higher value is a self-assessment further from \emph{Perfect}. Every response is a multiple of ten, so differences of a few points between cells are not interpreted.}
\label{tab:tlx_matrix} \\
\toprule
\textbf{Task} & \textbf{Mental} & \textbf{Physical} & \textbf{Temporal} & \textbf{Effort} & \textbf{Frustr.} & \textbf{Perf.} & \textbf{Aggregate} \\
\midrule
\endfirsthead
\toprule
\textbf{Task} & \textbf{Mental} & \textbf{Physical} & \textbf{Temporal} & \textbf{Effort} & \textbf{Frustr.} & \textbf{Perf.} & \textbf{Aggregate} \\
\midrule
\endhead
\midrule
\multicolumn{8}{r}{\textit{Continued on next page}}\\
\bottomrule
\endfoot
\bottomrule
\endlastfoot

3 Gherkin & 26.9 & 17.7 & 21.5 & 23.8 & 13.1 & 14.6 & 19.6 \\
4 Design & 47.7 & 21.5 & 40.8 & 36.2 & 20.0 & 16.2 & 30.4 \\
6 Packs & 40.0 & 21.5 & 32.3 & 36.2 & 26.2 & 18.5 & 29.1 \\
7 \ac{QA} & 25.4 & 13.8 & 19.2 & 23.1 & 15.4 & 10.8 & 18.0 \\
8 Deploy & 25.4 & 20.8 & 26.9 & 29.2 & 22.3 & 21.5 & 24.3 \\
9 Export & 9.2 & 7.7 & 6.9 & 7.7 & 3.1 & 10.8 & 7.6 \\ \addlinespace
\midrule
\textbf{All tasks} & \textbf{29.1} & \textbf{17.2} & \textbf{24.6} & \textbf{26.0} & \textbf{16.7} & \textbf{15.4} & \textbf{21.5} \\
\end{longtable}
}

Collapsed across all seventy-eight administrations the grand aggregate is 21.5 out of 100, and the profile has a definite shape. Mental demand is the largest component at 29.1, followed by effort at 26.0 and temporal demand at 24.6; frustration, at 16.7, is among the smallest, below physical demand. The burden this tool imposes therefore sits in comprehension rather than in struggle, which is the distinction the subscale profile exists to make and which a single aggregate would have flattened. Performance, reported as marked with its reversed anchors, averages 15.4, meaning participants placed their own success close to the \emph{Perfect} end throughout.

By aggregate, Design and Packs are the two heaviest tasks, the two in which the largest volume of generated content is put in front of the participant for review; Export, which generates nothing, is lightest by a wide margin. That the workload tracks the quantity of generated material to be assessed, rather than the number of interface steps, is consistent with the framework's own account of where it puts the practitioner's effort.

Three predictions were recorded before administration and two of them were wrong, which is reported here rather than left out. Task 3 was predicted to carry the session's highest mental demand; it carried the second highest, 26.9 against Design's 47.7. Task 6 was predicted to be moderate throughout on the grounds that it is largely a review task; it came second by aggregate, tied for the highest effort score at 36.2, and carried the highest frustration of any task at 26.2. Two predictions held: task 9 was predicted lowest overall and was lowest on every single subscale, and physical demand was predicted to sit near the floor and never exceeded 21.5. Task 8, predicted moderate throughout following its redesign away from a refusal test, came third by aggregate but carried the highest Performance value of any task at 21.5, meaning participants rated their own success on the deployment further from \emph{Perfect} than anywhere else in the session. \Cref{sec:eval_discussion} reads that one against an independently recorded defect in the same feature.

\subsection{Apex UX Questionnaire}
\label{sec:results_ux}

\Cref{tab:ux_items} reports all seventeen administered items individually, as response counts and a median, with no composite score and no construct means, for the reason given in \Cref{sec:apex_ux}. Responses of \emph{N/A - did not use this} are counted separately and excluded from the median.

{\footnotesize
\renewcommand{\arraystretch}{1.2}
\setlength{\tabcolsep}{4pt}
\begin{longtable}{l >{\raggedright\arraybackslash}p{8.2cm} c c c c c c c}
\caption[Apex \acs{UX} questionnaire responses, item by item]{Apex \acs{UX} questionnaire responses, item by item, $N = 13$. Scale 1 = strongly disagree to 5 = strongly agree; the \emph{N/A} column counts participants who marked \emph{did not use this}, which is excluded from the median. Items E2 and E3 evaluate the framework, not the instantiation.}
\label{tab:ux_items} \\
\toprule
& \textbf{Item} & \textbf{1} & \textbf{2} & \textbf{3} & \textbf{4} & \textbf{5} & \textbf{N/A} & \textbf{Mdn} \\
\midrule
\endfirsthead
\toprule
& \textbf{Item} & \textbf{1} & \textbf{2} & \textbf{3} & \textbf{4} & \textbf{5} & \textbf{N/A} & \textbf{Mdn} \\
\midrule
\endhead
\midrule
\multicolumn{9}{r}{\textit{Continued on next page}}\\
\bottomrule
\endfoot
\bottomrule
\endlastfoot

\multicolumn{9}{l}{\textit{A. Understanding where you are}} \\*
A1 & Knew which phase of the process I was in & 0 & 0 & 3 & 4 & 6 & 0 & 4 \\*
A2 & Understood what the tool expected next & 0 & 0 & 4 & 5 & 4 & 0 & 4 \\*
A3 & Could tell finished work from unfinished & 0 & 0 & 4 & 7 & 2 & 0 & 4 \\ \addlinespace
\multicolumn{9}{l}{\textit{B. Working with AI-generated content}} \\*
B1 & Could tell generated content from my own & 0 & 0 & 0 & 3 & 7 & 1 & 5 \\*
B2 & Reviewing took less effort than writing it & 0 & 0 & 2 & 7 & 4 & 0 & 4 \\*
B3 & Felt able to reject or change the output & 0 & 0 & 0 & 1 & 11 & 1 & 5 \\*
B4 & Trusted it enough to use in a real project & 0 & 0 & 3 & 8 & 2 & 0 & 4 \\*
B5 & Could see what information the AI was given & 0 & 1 & 2 & 6 & 4 & 0 & 4 \\ \addlinespace
\multicolumn{9}{l}{\textit{C. Control and reversibility}} \\*
C1 & Felt in control of what the tool did & 0 & 1 & 3 & 8 & 1 & 0 & 4 \\*
C2 & Confident I could undo or redo a step & 0 & 1 & 1 & 2 & 8 & 1 & 5 \\*
C3 & Understood the consequences of locking a phase & 0 & 1 & 4 & 2 & 6 & 0 & 4 \\*
C4 & Waiting for the AI to finish was acceptable & 0 & 0 & 0 & 4 & 9 & 0 & 5 \\ \addlinespace
\multicolumn{9}{l}{\textit{D. Feedback and errors}} \\*
D1 & Was told clearly what had happened when something went wrong & 0 & 1 & 6 & 4 & 2 & 0 & 3 \\*
D2 & Messages said what to do next, not only what failed & 0 & 0 & 4 & 3 & 6 & 0 & 4 \\ \addlinespace
\multicolumn{9}{l}{\textit{E. Fit to real work}} \\*
E1 & The artefacts are ones my team would use & 0 & 1 & 3 & 5 & 1 & 3 & 4 \\*
E2 & \textit{(framework)} The process fits how my team works & 0 & 2 & 2 & 6 & 0 & 3 & 4 \\*
E3 & \textit{(framework)} The effort is worth what it produces & 0 & 0 & 2 & 5 & 6 & 0 & 4 \\
\end{longtable}
}

Fifteen of the seventeen items are instantiation-scoped and are read first. Ten of the fifteen have a median of 4 and four have a median of 5: being able to tell generated content from one's own (B1), feeling able to reject or change it (B3), confidence that a step could be undone (C2), and the acceptability of waiting for generation to finish (C4). One item sits lower than every other closed item in the instrument: \textbf{D1}, on being told clearly what had happened when something went wrong, has a median of \textbf{3}, with six of the thirteen at the neutral point and one below it. Its paired item D2, on whether messages said what to do next rather than only what had failed, sits at 4, so what participants report is not an absence of messages but an absence of a legible account of what went wrong.

The trust-calibration pair predicted in \Cref{sec:apex_ux} to be the most diagnostic does not diverge at the median: B4, trust sufficient for real use, and B2, reviewing being cheaper than writing, both sit at 4, with no participant below 3 on either. On its face this is calibrated trust, which is the reassuring reading. It is qualified by the free text of a single participant who describes the opposite state directly, having believed the models were properly grounded and so, in their own words, ``I did not review in great detail and I trusted it completely'' (P7). One self-report is not a rate, and the design cannot say how many others were in the same position while answering 4, because nothing was observed. The limit on what that establishes is carried into \Cref{sec:threats}.

E3, on whether the effort of following the process is worth what it produces, is among the more strongly supported items in the instrument, with none disagreeing. E2, on whether the process fits the way the participant's team works, is the weakest-supported item of the seventeen. Three participants marked \emph{did not use this}, no one strongly agreed, two of the ten who did answer disagreed and two more were neutral, so its median of 4 rests on the six remaining agreements. The three \emph{N/A} responses on E2 come from the same three participants who marked \emph{N/A} on E1, so the missing values are not independent of one another; those three answered E3 normally. That a majority endorses the process as worth its cost while fewer will say it fits their own team is the most interesting result in this instrument, and it is a question about the framework rather than about Apex.

\subsection{Analytical Assessment}
\label{sec:results_analytical}

\Cref{tab:analytical} gives the rating for each of the five criteria stated in \Cref{sec:analytical}, with the evidence each rests on. No criterion is rated \emph{not met}, and that is not a favourable finding in itself: these are the framework's own design criteria, so a \emph{not met} would record a failure at something it explicitly set out to do. The information is in the three \emph{partially met} rows, each of which names a specific shortfall rather than a general reservation.

{\footnotesize
\renewcommand{\arraystretch}{1.2}
\setlength{\tabcolsep}{4pt}
\begin{longtable}{>{\raggedright\arraybackslash}p{2.5cm}
                  >{\raggedright\arraybackslash}p{1.9cm}
                  >{\raggedright\arraybackslash}p{10.7cm}}
\caption[Analytical assessment of the framework against criteria fixed in advance]{Analytical assessment of the framework against the five criteria fixed in \Cref{sec:analytical}, with the evidence each rating rests on.}
\label{tab:analytical} \\
\toprule
\textbf{Criterion} & \textbf{Rating} & \textbf{Justification and evidence} \\
\midrule
\endfirsthead
\toprule
\textbf{Criterion} & \textbf{Rating} & \textbf{Justification and evidence} \\
\midrule
\endhead
\midrule
\multicolumn{3}{r}{\textit{Continued on next page}}\\
\bottomrule
\endfoot
\bottomrule
\endlastfoot

Clarity & Partially met & Every construct in \Cref{tab:apex_mapping} resolved to a nameable mechanism, which is the operational test of whether a definition is precise enough to act on, and a practitioner described the central construct back unprompted, as a library for project information carrying guidelines, requirements, plans and tests together (T2, \Cref{sec:interviews}). Against that, spec-drift detection was found to presume a premise the framework never states, that the specification is scoped per story, and the omission surfaced only when the implementation failed on it (\Cref{sec:apex_limitations}). A definition whose preconditions are discoverable only by building it is not yet clear. \\ \addlinespace

Completeness & Partially met & Strategic Alignment and six lifecycle phases through to Maintenance are addressed, with three named loop-backs, so no stage of the lifecycle the framework claims is left uncovered (\Cref{tab:phase_mapping}, \Cref{fig:lifecycle}). Two of its own measurements fall short of their definitions: AI Defect Escape Rate is not instantiable at the observation point it presumes and exists only as a weaker proxy, and spec-drift detection holds for the functional specification but not the technical one (\Cref{tab:apex_mapping}). The gap is in the definitions rather than the coverage, and it is identifiable, which separates this rating from \emph{not met}. \\ \addlinespace

Alignment with \ac{SDLC} phases & Met & Each phase maps onto a recognised lifecycle stage with a defined AI contribution, an accountable hat and an exit control (\Cref{tab:phase_mapping}), and the decomposition survived instantiation with no phase merged, split or skipped (\Cref{tab:apex_mapping}). No interviewee questioned the phase set. One practitioner questioned its linearity (T4, \Cref{sec:interviews}), which is a claim about the transitions rather than the stages, and was answered by making the loop-backs out of Maintenance explicit rather than by altering the phases. \\ \addlinespace

Support for human-in-the-loop practice & Met & Every phase terminates in a write to the specification store and a human acceptance, so none can be exited by generation alone (\Cref{fig:apex_userflow,fig:apex_userflow_b}); the Interactive Bridge, the Consistency Factor and server-side reconciliation of generated claims are all realised (\Cref{tab:apex_mapping}). Participants report the preconditions for review independently: median 5 on telling generated content from their own, and eleven of thirteen strongly agreeing they could reject it (\Cref{tab:ux_items}). The criterion concerns whether review is structurally supported, not whether it happens; one participant reports not reviewing in detail at all (P7), and nothing in the framework would detect a gate satisfied without scrutiny. \\ \addlinespace

Governance of AI-enabled activity & Partially met & Gate enforcement and traceability are machine-checkable rather than procedural: per-story phase status refused by middleware, stable identifiers resolving every artefact to a requirement, and a Context Traceability Rate computed as a real aggregate (\Cref{tab:apex_mapping}). Two shortfalls are specific. The framework's own measure of AI-attributable defects exists only as a proxy over reported, linked, post-deployment issues, so the regime cannot see most of what it governs (\Cref{sec:apex_limitations}). And during the study a governance surface reported a verdict its evidence did not support, the Phase 5 pre-flight check reading a stale repository-context file as current and presenting it at high confidence, with four participants' task 8 free text converging on that surface (\Cref{sec:eval_discussion}). A control that can report confidently from stale evidence is not yet trustworthy without an independent check. \\

\end{longtable}
}

% #############################################################################
\section{Threats to Validity}
\label{sec:threats}

\paragraph{Construct validity.} Every instrument here measures perception, and none measures delivered software quality. \ac{SUS} and Raw \ac{NASA-TLX} are standardised and comparable against published corpora, which is the mitigation that does apply; the \ac{UX} questionnaire's construct-validity limit is that it is author-written and unvalidated, addressed by treating it as such (\Cref{sec:apex_ux}). Two departures already described (\Cref{sec:tlx}) reduce what the workload instrument can resolve: the 0 to 10 response scale, and the anchoring and fatigue risk of six repeated administrations. One respondent's data is consistent with exactly that: thirty-five of P6's thirty-six workload ratings are the same number. No mitigation was available for either and none is claimed.

The sharpest construct gap concerns phase gates: the item A4/task-8 cross-read described in \Cref{sec:apex_ux}, this study's only planned evidence on whether a gate reads as a guardrail or as obstruction, did not survive task 8's redesign away from a refusal test (\Cref{sec:tlx}). The gap is therefore worse than the absence of an observer check: there was no guaranteed refusal trigger for any of the thirteen participants, so gate comprehension is unmeasured rather than weakly measured. The framework's most distinctive mechanism is the one this evaluation has least evidence about, and that is stated rather than compensated for.

\paragraph{Internal validity.} The author designed both artefacts, wrote the task script and the seventeen \ac{UX} items, built the seeded demonstration project, and conducted the interviews. The unmoderated design removes that conflict from the sessions structurally (\Cref{sec:eval_instruments_procedure}), the strongest mitigation available, but removes nothing from the instruments themselves. Two incidents make the residue concrete rather than hypothetical. A configuration defect let an interface-language setting leak between projects, so a nominally English session was answered by the model in Portuguese, and it was noticed only because the author recognised the symptom in a participant's free text and root-caused it himself. A second defect, in which Phase 5 read a stale repository-context file as current and reported at high confidence, was found during the first session and was still present for every subsequent participant. In both cases the person best placed to detect and interpret the anomaly was the artefact's designer.

The same window produced a second internal threat: the system was not frozen. One defect fix was deployed during data collection and one configuration change followed it, so participants did not all use a byte-identical build. Against that, the project environment was reset to an identical seeded state before every session (\Cref{sec:tlx}).

\paragraph{External validity.} The sample (\Cref{sec:eval_participants}) is a convenience sample from one person's professional network, skewed junior, with no open call and no recorded count of those approached. It supports a description of how these thirteen people experienced the tool and no claim about experienced practitioners, teams or organisations. The pre-registered concern about a single organisational setting turns out to apply in a different place than expected: the participants are not from one organisation, but the \emph{task} environment is a single seeded project on one project-management tool, which ten of the thirteen had never used, so unfamiliarity with that board is confounded with unfamiliarity with Apex throughout. The observation window is one session of roughly an hour, which precludes any claim about long-term maintainability, adoption or organisational change; one participant makes the point himself, expecting that ``with regular use, I believe the user would become very familiar with the interface'' (P12), which is a hypothesis this design cannot test. The Portuguese figures rest on three respondents administered an instrument whose own validation reports weak inter-rater agreement, and are reported as indicative only.

\paragraph{Conclusion validity.} Thirteen participants spread over five levels of prior AI exposure cannot support inferential statistics or subgroup comparison, and none is computed. Two administration artefacts show what an unmoderated design cannot catch as it happens: one participant declared Portuguese as the session language and then completed every instrument in English, and another declared English but submitted the Portuguese consent form. Neither was correctable at the time because nobody was present to correct it, and both are reported as recorded.

\paragraph{The cost of recording nothing.} Two threats follow from the decision not to observe or record, and both are consequences of a deliberate design rather than oversights. First, there is no observational evidence of any kind: no measured completion rate, no time on task beyond a form-submission timestamp that cannot distinguish a hard task from a coffee break, no assist or error counts, and no independent check on comprehension anywhere. This leaves the perception-versus-quality threat standing rather than partly answered, and it is why the trust-calibration pair in \Cref{sec:results_ux} can report that trust and review effort did not diverge without being able to say whether any individual accepted output unread. Second, the interview analysis rests on contemporaneous written notes rather than transcripts (\Cref{sec:eval_participants}), and the selection was made by the person who designed the artefact under discussion; \Cref{sec:interviews} names which interviewee raised which point precisely so that the selection is at least attributable.

% #############################################################################
\section{Discussion}
\label{sec:eval_discussion}

The instruments do not tell one story, and the places they part company are more informative than the places they coincide.

The first apparent disagreement is between \ac{SUS} and the \ac{UX} questionnaire. A grade D mean of 56.73 sits alongside medians of 5 on being able to reject generated output, on confidence that a step could be undone, and on the acceptability of waiting for generation. That is only a contradiction if both are read as verdicts on the same object: \ac{SUS} asks about the system as a whole, three of its items specifically about complexity, about needing to learn a great deal and about needing technical support, while the \ac{UX} items ask about individual mechanisms. The free text resolves the two in the same direction: nine of the thirteen answers to what was most confusing name density, orientation or navigation rather than any AI mechanism, in terms such as ``too many things in the screen at the begging not sure where to click'' (P1) and ``too much happening at the same time'' (P6). The \ac{SUS} sub-scales point the same way, Learnability at 44.23 against Usability at 59.86. The reading all four sources support is that the mechanisms the framework requires are individually legible while the surface assembling them is not, and that the cost is concentrated at first contact. This is a finding about the instantiation and, for the reason given in \Cref{sec:eval_strategy}, it does not transfer to the framework.

The second cross-read is an agreement between instruments not designed to meet. Error and feedback surfacing is the weakest area the \ac{UX} questionnaire identifies, D1 having the lowest closed-item median. Independently, the two heaviest tasks by workload are Design and Packs, and Packs carries the highest frustration score of any task at 26.2 despite being predicted moderate on the grounds that it is largely review. A participant's free text supplies the mechanism connecting them, describing a mistake in one stage which ``Apex did not clearly show what the error was'' and then completing the following stages correctly but on that incorrect assumption (P8). A low D1 is not a cosmetic complaint about message wording: it describes a state in which a gate passes over an error that was never surfaced.

The third is the one place an instrument, a free-text pattern and an independently recorded defect converge. Task 8 was predicted to be moderate throughout after its redesign; it returned the highest Performance value of any task, 21.5, meaning participants rated their own success on the deployment furthest from \emph{Perfect}. Four participants' free text centres on the same surface, Phase 5's pre-flight reasoning about the repository, from three different angles: one read the ``infra changes required'' verdict as unexpected friction on a demonstration project, and two named the absence of a prior deployment as the resolving fact, which is the framework's own documented bootstrap rule rather than a fault. Recorded separately, and not by the instruments, is the defect described in \Cref{sec:threats}: that check read a stale context file as current and reported at high confidence. Three sources of evidence, one of which is not a questionnaire, point at the same mechanism, which is why \Cref{tab:analytical} rates governance \emph{partially met} rather than met.

Where the instruments genuinely disagree, they do so about trust, and the disagreement is not resolvable within this design. B4 and B2 both return a median of 4, which is the calibrated-trust signature the pair was built to detect. One participant's free text reports the uncalibrated state directly, having trusted completely and not reviewed in detail. The questionnaire cannot say whether that participant was alone or typical, because task 4's key observation was blanket acceptance of generated output and, nothing having been observed, no data on it exists. The honest position is that the instrument found no divergence and the study is not capable of finding one.

Finally, a disagreement running opposite to the one usually anticipated. The expected tension is a favourable usability score alongside interview scepticism about process overhead; here the tool score is below average while the framework's reception in the interviews is broadly favourable, including an unprompted endorsement of its central bet and an expert reading the phase and hat structure as a genuine departure from closed-role methodology (\Cref{sec:interviews}). A below-average \ac{SUS} is not evidence against the method, exactly as a good one would not have been evidence for it.

\subsection{Answering the Research Questions}
\label{sec:eval_rq_answers}

\paragraph{DRQ1, how can AI-supported activities be mapped to \ac{SDLC} phases so that the point, purpose and boundary of AI participation are unambiguous?} They can be, at phase granularity, and \Cref{tab:phase_mapping} is the answer: every phase names what AI may contribute, what it must produce, which hat is accountable, and the control that must be satisfied before exit. Two independent checks support it rather than the author's assertion. The decomposition survived instantiation with no phase merged, split or skipped, the alignment criterion rated \emph{met} in \Cref{tab:analytical}; and an expert who had not designed the framework read it back in interview, comparing its hats and its refusal to bind roles to individuals against the closed-role structures of organisational methodology (\Cref{sec:interviews}). The qualification is the clarity rating: boundaries are unambiguous at phase level, and at least one construct carried a precondition the framework never stated, which only building it exposed.

\paragraph{DRQ2, which governance mechanisms, quality gates and human-in-the-loop practices keep AI contributions traceable, validated and accountable?} Those named in \Cref{sec:governance,sec:gates}, and the evidence separates them by strength. Traceability and validation are the strongest: gates are enforced against specification content rather than against a signature, stable identifiers make every artefact resolve to a requirement, and generated claims about coverage are re-checked server-side because self-reporting was found untrustworthy in practice (\Cref{sec:apex_qa}). Human-in-the-loop support is rated \emph{met} on the structural test, no phase being exitable by generation alone, and participants confirm its preconditions, with a median of 5 on distinguishing generated content from their own and eleven of thirteen strongly agreeing they could reject it. Accountability is the weaker half and is rated \emph{partially met}: the framework's own measure of AI-attributable defects is instantiable only as a proxy, and a governance surface was observed reporting confidently on stale evidence. The mechanisms make contributions traceable and reviewable; they do not yet make the consequences of a contribution measurable.

\paragraph{DRQ3, to what extent can those constructs be operationalised in a working system, and which resist it?} Of the twelve constructs mapped in \Cref{tab:apex_mapping}, ten are realised as specified, one exists only as a named proxy, and one is partially realised with the resisting part identified. That count is the smaller part of the answer; the larger part is what the two exceptions revealed. Bolt Cycle Time was implemented, displayed and believed while measuring something else, and was caught only by comparing the implementation against the written definition, which establishes that a construct can be operationalised incorrectly without anything appearing wrong. Spec-drift detection failed because the framework assumes a per-story specification granularity it nowhere states, so the instantiation discovered an omission in the method rather than a defect in the tool. Both are reported as findings, and the existence of Apex on its own is not offered as an answer to this question.

\paragraph{DRQ4, does the framework, through that implementation, give practitioners usable and cognitively sustainable support in a real setting?} The honest answer is split, and the two halves do not contradict. On cognitive sustainability the evidence is favourable: a grand workload aggregate of 21.5 out of 100, frustration among the lowest subscales at 16.7, and a heaviest single task at 30.4. The burden is concentrated in mental demand and falls on the phases with the most generated content to review, which is where the framework intends to put it. On usability it is not favourable: a \ac{SUS} mean of 56.73, grade D and below the published corpus mean, with one participant at 15 and a Portuguese subgroup mean of 44.17 at grade F; error legibility is the weakest \ac{UX} item; and nine of thirteen named density or orientation as the most confusing thing about the tool. The two readings sit together because they measure different things: the tool asks little of the practitioner once they know where they are, and does too little to get them there. Both halves are claims about the instantiation, at one session of about an hour, with the sample limits stated in \Cref{sec:threats}; the framework's own component of DRQ4 rests on the demonstration in \Cref{chap:demonstration}.
